\subsection{Training pipeline}

The training pipeline is implemented in \texttt{train.py}, which integrates model optimization, hyperparameter configuration, checkpoint saving, and training monitoring via \texttt{wandb}.

We adopt a two-stage training strategy. In the first stage, we conduct exploratory experiments on a high-resolution setting (256$\times$256) with small datasize to quickly evaluate model behavior, verify training stability, and compare different loss functions and configurations. In this stage, a smaller batch size and a higher number of epochs are used to allow sufficient iteration over the limited data while keeping computational cost low. Specifically, the exploratory training is configured as:

\[
\text{optimizer} = \text{AdamW}, \quad
\text{lr} = 5 \times 10^{-4}, \quad
\text{weight decay} = 10^{-4}, \quad
\text{batch size} = 16, \quad
\text{epochs} = 50.
\]

Based on the observations and convergence behaviors from the exploratory stage, we proceed to the final training stage on the low-resolution setting (128$\times$128). The final configuration is selected to balance training efficiency and convergence stability. In particular, a larger batch size is adopted to stabilize gradient updates. The final training configuration is summarized as follows:

\[
\text{optimizer} = \text{AdamW}, \quad
\text{lr} = 5 \times 10^{-4}, \quad
\text{weight decay} = 10^{-4}, \quad
\text{batch size} = 64, \quad
\text{epochs} = 15.
\]